% Source: https://artofproblemsolving.com/community/c3h2294845_help_with_the_question
% Stable identity: aops-topic-2294845-q01
The proof is correct; the only step that needs justification is $h_1+h_2=h$.

Draw the line through $E$ perpendicular to $AB$. Since $CD\parallel AB$, it is also perpendicular to $CD$; let it meet line $AB$ at $F$ and line $CD$ at $G$. Then $EF=h_1$, $EG=h_2$ and $FG=h$.

\begin{asy}
usepackage("amsmath");
size(8cm);
defaultpen(fontsize(10pt));
pair A=(0,0), B=(6,0), C=(8,3.5), D=(2,3.5), X=(4.4,1.3), F=(4.4,0), G=(4.4,3.5);
fill(A--X--B--cycle,rgb(0.8,0.88,1)); fill(C--X--D--cycle,rgb(1,0.88,0.75));
draw(A--B--C--D--cycle,linewidth(0.9));
draw(A--X--B,blue); draw(C--X--D,orange);
draw(F--G,dashed);
real t=0.22;
draw((F.x,t)--(F.x+t,t)--(F.x+t,0)); draw((G.x,3.5-t)--(G.x+t,3.5-t)--(G.x+t,3.5));
dot(X); dot(F); dot(G);
label("$A$",A,SW); label("$B$",B,SE); label("$C$",C,NE); label("$D$",D,NW);
label("$E$",(3.97,1.36)); label("$F$",F,S); label("$G$",G,N);
label("$h_1$",(F.x,0.62),E); label("$h_2$",(F.x,2.45),E);
label("$x$",(1.6,0),S); label("$x$",(6.4,3.5),N);
\end{asy}

Because $E$ is inside the parallelogram, it lies in the strip between the parallel lines $AB$ and $CD$, so $E$ lies on the segment $FG$. Hence $FE+EG=FG$, that is, $h_1+h_2=h$.

The rest is the area formula for a triangle: triangle $AEB$ has base $AB=x$ and height $h_1$, and triangle $CED$ has base $CD=x$ and height $h_2$. So
\[[AEB]+[CED]=\dfrac12xh_1+\dfrac12xh_2=\dfrac12xh=\dfrac12[ABCD].\]
The student's argument is a complete proof. (The textbook proof, which draws the line through $E$ parallel to $AB$ and splits $ABCD$ into two parallelograms, each containing one of the triangles as half of it, proves the same thing with more steps.) $\blacksquare$
